\chapter{Processus à sauts}\label{ch:qm:jump-processes}

Le 19 octobre 1987, le Dow Jones Industrial Average a chuté de 22.6~\% en une séance. Mesurée à
l'aune de l'écart-type des rendements quotidiens du S\&P~500 sur six décennies, 0.98~\%, la chute de l'indice
ce jour-là représente un mouvement de 21 écarts-types. Un modèle gaussien avec cette volatilité donne à une telle
journée une probabilité d'environ $10^{-97}$~: il faudrait attendre plus longtemps que l'univers n'a
existé, et ce plusieurs fois. Aucune réestimation de la volatilité ne répare cela~; ce qui échoue, c'est
la continuité des trajectoires. Un processus qui évolue de manière continue, sauf à des instants aléatoires où il
saute, peut attribuer à cette même journée une probabilité d'une fois en cinquante ans, tout en laissant
les jours ordinaires restent ordinaires. Ce chapitre construit les processus à sauts utilisés dans la série~: \omterm{def:qm:jump-processes:poisson}{processus de Poisson} et
\omterm{def:qm:jump-processes:poisson}{processus de Poisson composé}, la formule d'Itô avec sauts et les \omterm{def:qm:jump-processes:levy}{processus de Lévy}, caractérisés par
une seule formule, la formule de Lévy--Khintchine, que les livres suivants utilisent pour évaluer les prix (le livre~5,
chapitre~13) et pour les transformer (chapitre~28 de ce livre).

\section{Processus de Poisson et processus de Poisson composé}

\begin{definition}[Processus de Poisson, processus de Poisson composé]\label{def:qm:jump-processes:poisson}
Un \emph{processus de Poisson}\index{processus de Poisson} de taux $\lambda > 0$ est un processus de comptage $N_t$ tel que $N_0 =
0$, à accroissements indépendants et $N_t - N_s \sim \mathrm{Poisson}(\lambda(t - s))$. Un \emph{processus de Poisson
composé}\index{processus de Poisson composé} est $J_t = \sum_{i=1}^{N_t}Y_i$, les amplitudes des sauts $Y_i$ étant
indépendantes, identiquement distribuées et indépendantes de $N$.
\end{definition}

Les délais entre événements sont des lois exponentielles indépendantes de moyenne $1/\lambda$, et le premier
instant d'arrivée $\tau_1$ vérifie $\P(\tau_1 > t) = e^{-\lambda t}$~: la probabilité de survie du livre~2,
chapitre~23, où le taux est une intensité de défaut. Un \omterm{def:qm:jump-processes:poisson}{processus de Poisson} n'est pas une \omterm{def:qm:probability-at-speed:martingale}{martingale}, mais il
devient une \omterm{def:qm:probability-at-speed:martingale}{martingale} lorsqu'on en retire la croissance espérée.

\begin{definition}[Processus de Poisson compensé]\label{def:qm:jump-processes:compensated}
Le \emph{processus de Poisson compensé}\index{processus de Poisson compensé} est $\tilde N_t = N_t - \lambda
t$~; pour un \omterm{def:qm:jump-processes:poisson}{processus de Poisson composé}, $J_t - \lambda\E[Y]t$.
\end{definition}

Les deux sont des \omterm{def:qm:probability-at-speed:martingale}{martingales}, avec $\E[\tilde N_t^2] = \lambda t$~; les accroissements sont indépendants et centrés,
et la variance d'une variable de Poisson est égale à sa moyenne. Le processus compensé est l'analogue à sauts
du \omterm{def:qm:brownian-motion:bm}{mouvement brownien}~: une \omterm{def:qm:probability-at-speed:martingale}{martingale} de \omterm{def:qm:brownian-motion:qv}{variation quadratique} $[\tilde N]_t = N_t$, le nombre de sauts
(la \cref{fig:qm:jump-processes:paths}, à gauche). Le chapitre~7 généralise le compensateur aux taux qui
dépendent du passé.

\section{Formule d'Itô avec sauts}

\begin{definition}[Diffusion à sauts, semimartingale]\label{def:qm:jump-processes:jd}
Une \emph{diffusion à sauts}\index{diffusion à sauts} est $X_t = X_0 + \int_0^t\mu_s\,ds + \int_0^t\sigma_s\,dW_s + J_t$,
où $J$ est un \omterm{def:qm:jump-processes:poisson}{processus de Poisson composé} (un nombre fini de sauts sur tout intervalle borné). Une
\emph{semimartingale}\index{semimartingale} est un processus somme d'une \omterm{def:qm:ito-calculus:local}{martingale locale} et d'un \omterm{def:qm:probability-at-speed:filtration}{processus adapté}
à variation finie, tous deux càdlàg~: la classe la plus large d'intégrateurs pour lesquelles les intégrales stochastiques
et la formule d'Itô fonctionnent.
\end{definition}

\begin{theorem}[Formule d'Itô avec sauts]\label{thm:qm:jump-processes:ito}
Pour une \omterm{def:qm:jump-processes:jd}{diffusion à sauts} $X$ et $f \in C^2$,
\[
f(X_t) = f(X_0) + \int_0^tf'(X_{s-})(\mu_s\,ds + \sigma_s\,dW_s) + \tfrac12\int_0^tf^{\prime\prime}(X_{s-})\sigma_s^2\,ds +
\sum_{s \le t}\bigl(f(X_s) - f(X_{s-})\bigr).
\]
\end{theorem}

\begin{proof}
Entre deux instants de saut, le processus est un \omterm{def:qm:ito-calculus:process}{processus d'Itô} et le \cref{thm:qm:ito-calculus:ito} s'applique~; à
un instant de saut $\tau_i$, la variation de $f(X)$ vaut $f(X_{\tau_i}) - f(X_{\tau_i -})$. Il y a un nombre fini de sauts
sur $[0, t]$, donc la somme des morceaux donne la formule.
\end{proof}

Le terme de saut est la variation exacte, et non un développement de Taylor~: les sauts ne sont pas petits. Avec $f = \exp$ et
$S = e^X$, $dS_t/S_{t-} = (\mu_t + \tfrac12\sigma_t^2)\,dt + \sigma_t\,dW_t + (e^{\Delta X_t} - 1)$~; en compensant le
dernier terme, $S$ est une \omterm{def:qm:probability-at-speed:martingale}{martingale} si et seulement si $\mu + \tfrac12\sigma^2 + \lambda(\E[e^Y] - 1) = 0$, la correction de dérive
que porte chaque modèle à sauts du livre~5.

\begin{omfigure}
\begin{tikzpicture}
\begin{groupplot}[group style={group size=2 by 1,horizontal sep=1.6cm},
  width=7.6cm,height=4.8cm,tick label style={font=\small},label style={font=\small},
  grid=major,grid style={gray!25},table/col sep=comma]
\nextgroupplot[xlabel={$t$ (années)},xmin=0,xmax=2,ymin=-2,ymax=12,
  legend style={font=\footnotesize,at={(0.03,0.97)},anchor=north west,fill=white,draw=none}]
\addplot[omDef,thick,const plot] table[x=t,y=n]{figdata/methods/06-jump-processes/poisson.csv};
\addplot[omThm,thick] table[x=t,y=comp]{figdata/methods/06-jump-processes/poisson.csv};
\legend{$N_t$,$N_t - 5t$}
\nextgroupplot[xlabel={$t$ (années)},ylabel={log-prix},xmin=0,xmax=10,ymin=-0.6,ymax=0.3]
\addplot[omDef,thin] table[x=t,y=x]{figdata/methods/06-jump-processes/jd_path.csv};
\addplot[omThm,only marks,mark=o,mark size=2.6pt] table[x=t,y=x]{figdata/methods/06-jump-processes/jd_jumps.csv};
\end{groupplot}
\end{tikzpicture}

\omcaption{À gauche~: un \omterm{def:qm:jump-processes:poisson}{processus de Poisson} de taux 5 par an et sa \omterm{def:qm:probability-at-speed:martingale}{martingale} compensée $N_t - 5t$.
À droite~: dix années de la \omterm{def:qm:jump-processes:jd}{diffusion à sauts} calibrée du chapitre (pas quotidiens~; sauts au taux de 0.5 par an
de moyenne $-11.1\%$, signalés), avec la dérive qui annule le rendement quotidien moyen. Données~: le
tutoriel du chapitre, graine fixée.}\label{fig:qm:jump-processes:paths}
\end{omfigure}

\section{Processus de Lévy et formule de Lévy--Khintchine}

\begin{definition}[Processus de Lévy, loi indéfiniment divisible]\label{def:qm:jump-processes:levy}
Un \emph{processus de Lévy}\index{processus de Lévy} est un processus càdlàg tel que $X_0 = 0$, à accroissements indépendants et
stationnaires. Une loi $\mu$ est une \emph{loi indéfiniment divisible}\index{loi indéfiniment divisible}
si, pour tout $n$, elle est la loi d'une somme de $n$ variables indépendantes et identiquement
distribuées.
\end{definition}

La loi de $X_1$ d'un \omterm{def:qm:jump-processes:levy}{processus de Lévy} est indéfiniment divisible ($X_1$ est la somme de $n$ accroissements de longueur
$1/n$), et réciproquement toute \omterm{def:qm:jump-processes:levy}{loi indéfiniment divisible} est la loi de $X_1$ pour un certain \omterm{def:qm:jump-processes:levy}{processus de Lévy}. Le
\omterm{def:qm:brownian-motion:drift}{mouvement brownien avec dérive}, les \omterm{def:qm:jump-processes:poisson}{processus de Poisson} et de Poisson composé, et leurs sommes sont des \omterm{def:qm:jump-processes:levy}{processus de Lévy}~; le
cas général ajoute une infinité de petits sauts.

\begin{definition}[Mesure de Lévy, triplet de Lévy, exposant caractéristique]\label{def:qm:jump-processes:triplet}
La \emph{mesure de Lévy}\index{mesure de Lévy} d'un \omterm{def:qm:jump-processes:levy}{processus de Lévy} est $\nu(A) = \E[\#\{s \le 1 : \Delta X_s \in A\}]$,
le nombre moyen de sauts par unité de temps dont la taille appartient à $A$ ($0 \notin \bar A$)~; elle vérifie
$\int\min(1, x^2)\,\nu(dx) < \infty$. Le \emph{triplet de Lévy}\index{triplet de Lévy} est $(\sigma^2, \nu, \gamma)$~:
variance gaussienne, mesure de Lévy et dérive. L'\emph{exposant caractéristique}\index{exposant caractéristique}
est la fonction $\psi$ telle que $\varphi_{X_t}(u) = e^{t\psi(u)}$.
\end{definition}

\begin{theorem}[Lévy--Khintchine]\label{thm:qm:jump-processes:lk}
L'\omterm{def:qm:jump-processes:triplet}{exposant caractéristique} d'un \omterm{def:qm:jump-processes:levy}{processus de Lévy} de triplet $(\sigma^2, \nu, \gamma)$ est
\[
\psi(u) = \iu\gamma u - \tfrac12\sigma^2u^2 + \int_{\R}\bigl(e^{\iu ux} - 1 - \iu ux\mathbf 1_{\{|x|<1\}}\bigr)\,\nu(dx),
\]
et tout triplet vérifiant $\sigma^2 \ge 0$ et $\int\min(1, x^2)\,\nu(dx) < \infty$ définit un \omterm{def:qm:jump-processes:levy}{processus de Lévy}.
De plus, $X$ se décompose en un \omterm{def:qm:brownian-motion:drift}{mouvement brownien avec dérive}, un \omterm{def:qm:jump-processes:poisson}{processus de Poisson composé} des sauts
de taille supérieure à un, et une \omterm{def:qm:probability-at-speed:martingale}{martingale} de carré intégrable des petits sauts compensés (Lévy--Itô).
\end{theorem}

\admitted

Pour un \omterm{def:qm:jump-processes:poisson}{processus de Poisson composé} d'intensité $\lambda$ et de loi des sauts $F$, $\nu = \lambda F$ et l'intégrale vaut
$\lambda(\varphi_Y(u) - 1)$, à une dérive près~: le calcul tient en une ligne par conditionnement sur $N_t$.
La \omterm{def:qm:jump-processes:jd}{diffusion à sauts} du chapitre, à sauts gaussiens, a pour exposant $\psi(u) = \iu\mu u - \tfrac12\sigma^2u^2 + \lambda(e^{\iu u\mu_J - \sigma_J^2u^2/2} -
1)$. Chaque modèle du projet en cours est spécifié par son exposant (le premier
listing du tutoriel), et la \cref{fig:qm:jump-processes:densities} compare quatre d'entre eux à l'horizon d'un jour.

\begin{definition}[Subordinateur]\label{def:qm:jump-processes:subordinator}
Un \emph{subordinateur}\index{subordinateur} est un \omterm{def:qm:jump-processes:levy}{processus de Lévy} à trajectoires croissantes~; en faisant courir
un \omterm{def:qm:brownian-motion:bm}{mouvement brownien} sur lui comme horloge, $X_t = \theta G_t + \sigma W_{G_t}$, donne un \omterm{def:qm:jump-processes:levy}{processus de Lévy} dont l'exposant est
l'exposant de Laplace du subordinateur évalué en $\iu u\theta - \sigma^2u^2/2$.
\end{definition}

Le \omterm{def:qm:jump-processes:subordinator}{subordinateur} gamma donne le processus variance gamma, et le \omterm{def:qm:jump-processes:subordinator}{subordinateur} gaussien inverse le processus normal
inverse Gaussian~: un temps de négociation qui s'écoule vite les jours chargés et lentement les jours calmes. Le
livre~5, chapitre~13, définit les modèles de valorisation bâtis sur eux. Avec la même variance annuelle que le
\omterm{def:qm:jump-processes:jd}{diffusion à sauts} du chapitre, elles répartissent autrement leur queue~: la \omterm{def:qm:jump-processes:jd}{diffusion à sauts} conserve un corps gaussien
et y ajoute une bosse lointaine de jours de krach, les modèles subordonnés épaississent régulièrement toute la queue
(la \cref{fig:qm:jump-processes:densities}).

\begin{omfigure}
\begin{tikzpicture}
\begin{axis}[width=11cm,height=5.2cm,ymode=log,
  xlabel={rendement logarithmique d'un jour (\%)},ylabel={densité},
  tick label style={font=\small},label style={font=\small},
  xmin=-25,xmax=8,ymin=1e-4,ymax=3,grid=major,grid style={gray!25},unbounded coords=jump,
  legend columns=4,legend style={font=\footnotesize,at={(0.5,-0.3)},anchor=north,draw=none,fill=none,/tikz/every even column/.append style={column sep=6pt}},
  table/col sep=comma]
\addplot[black,thick,dashed] table[x=x,y=gauss]{figdata/methods/06-jump-processes/densities.csv};
\addplot[omThm,thick] table[x=x,y=jd]{figdata/methods/06-jump-processes/densities.csv};
\addplot[omDef,thick] table[x=x,y=vg]{figdata/methods/06-jump-processes/densities.csv};
\addplot[omProp,thick] table[x=x,y=nig]{figdata/methods/06-jump-processes/densities.csv};
\legend{gaussien,diffusion à sauts,variance gamma,normale inverse gaussienne}
\end{axis}
\end{tikzpicture}

\omcaption{Densités du rendement logarithmique d'un jour, en échelle logarithmique, pour quatre modèles de Lévy de même
variance (écart-type quotidien de 0.98~\%)~: le modèle gaussien, la \omterm{def:qm:jump-processes:jd}{diffusion à sauts} calibrée du chapitre,
variance gamma ($\theta = -0.1$, $\nu = 0.013$) et normal inverse Gaussian ($\alpha = 58$, $\beta = -11.6$), les deux derniers avec des kurtosis d'excès à un jour de 8.5 et 12.5. Histogrammes de
2 millions de jours simulés chacun, avec des classes de largeur 0.5 point. Données~: le tutoriel du chapitre,
à graine fixée.}\label{fig:qm:jump-processes:densities}
\end{omfigure}

\section{Cumulants et structure par terme du kurtosis}

\begin{definition}[Cumulant]\label{def:qm:jump-processes:cumulant}
Le $n$-ième \emph{cumulant}\index{cumulant} $\kappa_n$ d'une variable aléatoire est la dérivée $n$-ième en
zéro de sa fonction génératrice des cumulants $\ln\E[e^{sX}]$~; $\kappa_1$ est la moyenne, $\kappa_2$ la
variance, $\kappa_3/\kappa_2^{3/2}$ l'asymétrie et $\kappa_4/\kappa_2^2$ le kurtosis d'excès.
\end{definition}

\begin{proposition}[Les cumulants d'un processus de Lévy sont proportionnels au temps]\label{prop:qm:jump-processes:scaling}
Pour un \omterm{def:qm:jump-processes:levy}{processus de Lévy}, $\kappa_n(X_t) = t\,\kappa_n(X_1)$. Ainsi l'asymétrie de $X_t$ décroît comme
$t^{-1/2}$ et son kurtosis d'excès comme $t^{-1}$. Pour des sauts gaussiens, $\kappa_n(X_1)$ est la partie
gaussienne plus $\lambda\E[Y^n]$.
\end{proposition}

\begin{proof}
$\ln\E[e^{sX_t}] = t\psi(-\iu s)$ est linéaire en $t$, de même que ses dérivées. Pour la partie
poissonienne composée, la fonction génératrice des \omterm{def:qm:jump-processes:cumulant}{cumulants} est $\lambda(\E[e^{sY}] - 1)$, dont la dérivée $n$-ième en zéro est
$\lambda\E[Y^n]$.
\end{proof}

Un modèle à sauts qui explique les krachs d'un jour est proche d'un modèle gaussien à un mois~: le kurtosis dont il a besoin à
un jour est dilué par la moyenne. Pour la calibration du chapitre, le kurtosis d'excès vaut 76 à un jour
et 3.6 à un mois (la \cref{fig:qm:jump-processes:kurtosis}). Les rendements mesurés sur les marchés conservent davantage
de kurtosis à un mois que ne le permet cette loi en $1/t$, car leur volatilité elle-même se regroupe~: c'est
l'indépendance des accroissements, et non les sauts, qui échoue aux horizons plus longs (chapitre~18).

\begin{omfigure}
\begin{tikzpicture}
\begin{groupplot}[group style={group size=2 by 1,horizontal sep=1.7cm},
  width=7.6cm,height=4.9cm,tick label style={font=\small},label style={font=\small},
  grid=major,grid style={gray!25},table/col sep=comma]
\nextgroupplot[xlabel={baisse d'un jour (\%)},ylabel={$\log_{10}\P(\text{baisse} \ge x)$},xmin=0,xmax=25,ymin=-100,ymax=5,
  legend style={font=\footnotesize,at={(0.03,0.03)},anchor=south west,fill=white,draw=none}]
\addplot[omDef,very thick] table[x=drop,y=normal]{figdata/methods/06-jump-processes/tail.csv};
\addplot[omThm,very thick] table[x=drop,y=jd]{figdata/methods/06-jump-processes/tail.csv};
\addplot[gray,dashed,forget plot] coordinates {(20,-100) (20,5)};
\legend{Gaussien,diffusion à sauts}
\nextgroupplot[xlabel={horizon (jours)},ylabel={kurtosis d'excès},xmode=log,ymode=log,xmin=0.8,xmax=320,ymin=0.2,ymax=120,
  legend style={font=\footnotesize,at={(0.97,0.97)},anchor=north east,fill=white,draw=none}]
\addplot[omDef,very thick] table[x=h,y=theory]{figdata/methods/06-jump-processes/kurtosis.csv};
\addplot[omThm,only marks,mark=*,mark size=1.7pt] table[x=h,y=sim]{figdata/methods/06-jump-processes/kurtosis.csv};
\legend{$\kappa_4/(t\,\kappa_2^2)$,simulé}
\end{groupplot}
\end{tikzpicture}

\omcaption{À gauche~: probabilité que le rendement d'un jour baisse d'au moins $x$, en échelle $\log_{10}$, pour une
loi gaussienne de l'écart-type quotidien du marché (0.98~\%) et pour la \omterm{def:qm:jump-processes:jd}{diffusion à sauts} du chapitre, de même
variance~; à 20~\% (en pointillés), elles diffèrent de 88 ordres de grandeur. À droite~: le kurtosis
d'excès de la \omterm{def:qm:jump-processes:jd}{diffusion à sauts} sur des horizons de 1 à 252 jours, à partir de ses \omterm{def:qm:jump-processes:cumulant}{cumulants} et de 400\,000 rendements simulés par
horizon. Données~: le tutoriel du chapitre, à graine fixée.}\label{fig:qm:jump-processes:kurtosis}
\end{omfigure}

\section{Tutoriel~: une fonction caractéristique vérifiable}\label{tut:qm:jump-processes:charfn}

\begin{tutorial}
\textbf{Objectif.} Spécifier la \omterm{def:qm:jump-processes:jd}{diffusion à sauts} du chapitre par son \omterm{def:qm:jump-processes:triplet}{exposant caractéristique}, la simuler et
vérifier l'exposant, les \omterm{def:qm:jump-processes:cumulant}{cumulants} et la structure par terme du kurtosis à l'aide des trajectoires.
\textbf{État final~:} les \cref{fig:qm:jump-processes:paths,fig:qm:jump-processes:kurtosis,fig:qm:jump-processes:charfn}
et la calibration du problème du week-end.
\begin{tutsteps}
\item \textbf{Les exposants} des modèles de Lévy du projet en cours.
\omcode{code/firm/levy/firm_levy.py}{26}{48}{Exposants caractéristiques~: mouvement brownien, sauts gaussiens et sauts double-exponentiels, variance gamma, normal inverse Gaussian.}\label{lst:qm:jump-processes:psi}
\item \textbf{La calibration}~: procéder par dichotomie sur la taille moyenne des sauts jusqu'à ce qu'une baisse de 20~\% soit un
événement cinquantennal, en maintenant la variance quotidienne à $0.98\%^2$.
\omcode{code/methods/06-jump-processes/python/qm_jumps.py}{45}{61}{Calibration de la taille des sauts sur une fréquence de queue à variance fixée.}\label{lst:qm:jump-processes:calibrate}
\item \textbf{Exécuter} \texttt{problem()}, \texttt{charfn\_check()} et \texttt{fig\_jumps.py}~; les \omterm{def:qm:jump-processes:cumulant}{cumulants} numériques de
\texttt{firm\_levy.cumulants} concordent avec les formes closes à $10^{-7}$.
\end{tutsteps}
\textbf{Pour aller plus loin.} Remplacer les sauts gaussiens par des sauts double-exponentiels de même moyenne et
variance, puis comparer les deux queues à 20~\%~; simuler le processus variance gamma par subordination et vérifier son
kurtosis par rapport à $3\nu$ à temps unitaire.
\end{tutorial}

\begin{omfigure}
\begin{tikzpicture}
\begin{axis}[width=11cm,height=4.6cm,
  xlabel={$u$},ylabel={$\varphi_{X_t}(u)$},
  tick label style={font=\small},label style={font=\small},
  xmin=0,xmax=40,ymin=-0.1,ymax=1.05,grid=major,grid style={gray!25},
  legend columns=2,legend style={font=\footnotesize,at={(0.5,-0.3)},anchor=north,draw=none,fill=none,/tikz/every even column/.append style={column sep=8pt}},
  table/col sep=comma]
\addplot[omDef,very thick] table[x=u,y=re_th]{figdata/methods/06-jump-processes/charfn.csv};
\addplot[omDef,only marks,mark=*,mark size=1.4pt] table[x=u,y=re_emp]{figdata/methods/06-jump-processes/charfn.csv};
\addplot[omThm,very thick] table[x=u,y=im_th]{figdata/methods/06-jump-processes/charfn.csv};
\addplot[omThm,only marks,mark=*,mark size=1.4pt] table[x=u,y=im_emp]{figdata/methods/06-jump-processes/charfn.csv};
\legend{partie réelle~: $e^{t\psi(u)}$,partie réelle~: simulée,partie imaginaire~: $e^{t\psi(u)}$,partie imaginaire~: simulée}
\end{axis}
\end{tikzpicture}

\omcaption{\omterm{def:qm:probability-at-speed:charfn}{Fonction caractéristique} du rendement à un mois du processus à sauts calibré~: la
formule de Lévy--Khintchine (traits) face à la moyenne empirique de $e^{\iu uX_t}$ sur 200\,000
rendements simulés (points). Données~: le tutoriel du chapitre, avec graine fixée.}\label{fig:qm:jump-processes:charfn}
\end{omfigure}

\section{Construction~: la boîte à outils de Lévy}\label{bld:qm:jump-processes:levy}

\begin{build}
\bfield{Rôle} Chaque modèle de sauts de la mini-entreprise est spécifié une seule fois, par son \omterm{def:qm:jump-processes:triplet}{exposant
caractéristique}, et utilisé de trois façons~: pour simuler, pour calculer les moments et pour évaluer par
transformée (chapitre~28).
\bfield{Interface} \texttt{psi\_bm}, \texttt{psi\_merton}, \texttt{psi\_kou}, \texttt{psi\_vg},
\texttt{psi\_nig}~; \texttt{cumulants(psi, n\_max)}~; \texttt{merton\_cumulants(\dots)}~;
\texttt{exp\_martingale\_drift(psi)}~; \texttt{simulate\_merton}, \texttt{simulate\_vg},
\texttt{simulate\_nig}.
\bfield{Règles} Les exposants prennent des arguments complexes et suivent la convention de la série $\varphi_{X_t}(u) =
e^{t\psi(u)}$~; les simulations sont exactes sur leur grille~; la dérive de \omterm{def:qm:probability-at-speed:martingale}{martingale} est calculée à partir de l'exposant,
jamais saisie à la main.
\bfield{Tests d'acceptation} \texttt{code/firm/levy/tests/}~: les \omterm{def:qm:jump-processes:cumulant}{cumulants} numériques sont égaux aux formes closes~;
les \omterm{def:qm:probability-at-speed:charfn}{fonctions caractéristiques} empiriques des trajectoires simulées concordent avec $e^{t\psi(u)}$~; les trajectoires
variance gamma et gaussienne inverse normale ont les moments que leurs exposants donnent~; $e^X$ avec la dérive
calculée a une moyenne constante.
\bfield{Pour aller plus loin} Exposants CGMY et stables~; simulation exacte des \omterm{def:qm:brownian-motion:passage}{temps de premier passage} pour les
processus spectralement négatifs.
\end{build}

\begin{omsources}
\item Federal Reserve History, «~Stock Market Crash of 1987~».
\item CFA Institute, Enterprising Investor, «~Fact file~: S\&P~500 sigma events~», 2012.
\item R.~C.~Merton, «~Option pricing when underlying stock returns are discontinuous~»,
\emph{Journal of Financial Economics} 3, 1976.
\item S.~G.~Kou, «~A jump-diffusion model for option pricing~», \emph{Management Science} 48, 2002.
\item R.~Cont and P.~Tankov, \emph{Financial Modelling with Jump Processes}, Chapman \& Hall/CRC, 2004.
\end{omsources}

\section{Exercices}

\begin{exercise}[$\star$]\label{exo:qm:jump-processes:1}
Les événements arrivent selon un \omterm{def:qm:jump-processes:poisson}{processus de Poisson} à raison de 5 par an. Calculer la probabilité qu'il n'y en ait aucun en un mois, et le
temps d'attente moyen avant le suivant.
\end{exercise}

\begin{exercise}[$\star$]\label{exo:qm:jump-processes:2}
Les sauts arrivent à raison de 3 par an, avec des tailles $\mathcal N(-2\%, 4\%^2)$. Calculer la moyenne et l'écart-type de
leur somme sur un an.
\end{exercise}

\begin{exercise}[$\star$]\label{exo:qm:jump-processes:3}
Montrer que $N_t - \lambda t$ est une \omterm{def:qm:probability-at-speed:martingale}{martingale} et que $\E[(N_t - \lambda t)^2] = \lambda t$.
\end{exercise}

\begin{exercise}[$\star\star$]\label{exo:qm:jump-processes:4}
Pour $X = \sigma W + J$ avec $J$ composé de Poisson, écrire $d(e^{X_t})$ à l'aide du \cref{thm:qm:jump-processes:ito}, et
déterminer la dérive à ajouter à $X$ pour que $e^X$ soit une \omterm{def:qm:probability-at-speed:martingale}{martingale}.
\end{exercise}

\begin{exercise}[$\star\star$]\label{exo:qm:jump-processes:5}
Les sauts arrivent à un taux de 2, avec des tailles exponentielles de moyenne $1/50$. Écrire l'\omterm{def:qm:jump-processes:triplet}{exposant caractéristique} du
\omterm{def:qm:jump-processes:poisson}{processus de Poisson composé}, puis sa moyenne et sa variance en $t = 1$.
\end{exercise}

\begin{exercise}[$\star\star$]\label{exo:qm:jump-processes:6}
Un modèle de Lévy présente un excès de kurtosis de 76.4 à un jour. Calculer l'excès de kurtosis à une semaine de cinq jours, puis
à une année de 252 jours.
\end{exercise}

\begin{exercise}[$\star\star\star$]\label{exo:qm:jump-processes:7}
\emph{Codage.} À l'aide de \texttt{cumulants}, calculer les quatre premiers \omterm{def:qm:jump-processes:cumulant}{cumulants} du processus variance gamma
avec $\theta = -0.1$, $\sigma = 0.2$, $\nu = 0.3$ en temps unitaire, et les comparer à $\kappa_1 = \theta$, $\kappa_2 = \sigma^2 +
\nu\theta^2$, $\kappa_3 = 2\theta^3\nu^2 + 3\sigma^2\theta\nu$, $\kappa_4 = 3\sigma^4\nu + 6\theta^4\nu^3 + 12\sigma^2\theta^2\nu^2$.
\end{exercise}

\begin{exercise}[$\star\star\star$]\label{exo:qm:jump-processes:8}
\emph{Trouver l'erreur.} «~Les rendements mensuels de notre indice présentent un excès de kurtosis de 1.2, proche de la loi gaussienne,
donc les sauts ne comptent pas pour notre risque à un jour.~»
\end{exercise}

\section{Problème~: vingt écarts-types}

\begin{problem}[{Problème du week-end --- une chute qu'une diffusion ne peut pas produire}]\label{pb:qm:jump-processes:1}
Les rendements quotidiens d'un indice ont un écart-type de 0.98~\%. Un gestionnaire des risques veut un modèle dans lequel une
chute d'un jour de 20~\% ou plus survient une fois en cinquante ans de 252 jours, et où les jours ordinaires conservent cette
variance. Le modèle est une \omterm{def:qm:jump-processes:jd}{diffusion à sauts}~: des sauts d'intensité $\lambda = 0.5$ par an, de tailles $\mathcal N(\mu_J,
0.05^2)$, et une diffusion dont la volatilité représente le reste de la variance.

\medskip\noindent\textbf{Partie I --- Le jour.}
\begin{enumerate}
\item Quelle est l'ampleur d'une chute de 20.98 écarts-types~?
\item Quelle probabilité une loi gaussienne de même écart-type lui attribue-t-elle~? (Donner $\log_{10}$.)
\item Et pour une chute d'exactement 20~\%~?
\item Quelle volatilité annuelle correspond à 0.98~\% par jour~?
\item Pourquoi réestimer la volatilité ne suffit-il pas à sauver le modèle gaussien~?
\end{enumerate}

\medskip\noindent\textbf{Partie II --- La \omterm{def:qm:jump-processes:jd}{diffusion à sauts}.}
\begin{enumerate}[resume]
\item Quelle taille moyenne de saut $\mu_J$ fait d'une chute de 20~\% un événement qui ne survient qu'une fois tous les cinquante ans~?
\item Quelle volatilité de diffusion maintient la variance quotidienne à $0.98\%^2$~?
\item Quelle part de la variance les sauts portent-ils~?
\item Quelle est la probabilité d'un saut un jour donné~?
\item À quelle fréquence le modèle produit-il une chute de 10~\% ou plus~?
\end{enumerate}

\medskip\noindent\textbf{Partie III --- \omterm{def:qm:jump-processes:cumulant}{Cumulants}.}
\begin{enumerate}[resume]
\item Quels sont l'asymétrie et le kurtosis d'excès à un jour~?
\item Quel est le kurtosis d'excès pour un mois de 21 jours~?
\item Pourquoi décroît-il en $1/t$~?
\item Quelle correction de dérive fait de $e^X$ une \omterm{def:qm:probability-at-speed:martingale}{martingale}~?
\item Vérifier le kurtosis à cinq jours par simulation.
\end{enumerate}

\medskip\noindent\textbf{Partie IV --- Jugement.}
\begin{enumerate}[resume]
\item « Une fois tous les cinquante ans » est-il une estimation ou un choix~?
\item Que manquerait le modèle que les données montrent à un mois~?
\item Pour quelles options la calibration importe-t-elle le plus~?
\item Énoncer le \emph{résultat nommé}~: l'ampleur de la journée de 1987 sous une diffusion, et le saut qui la rend rare mais possible.
\item En une phrase~: que change un processus à sauts que ne peut pas changer une volatilité~?
\end{enumerate}
\end{problem}

\section{Questions d'entretien}

\begin{interviewq}[$\star$ \iqroles{trader, risque}]\label{iq:qm:jump-processes:1}
Un modèle affirme que le mouvement d'hier était un événement à 20 écarts-types. Que concluez-vous~?
\end{interviewq}

\begin{interviewq}[$\star$ \iqroles{chercheur}]\label{iq:qm:jump-processes:2}
Les transactions arrivent selon un \omterm{def:qm:jump-processes:poisson}{processus de Poisson} à 2 par seconde. Quelle est la probabilité d'en avoir au moins deux en une seconde, et
quelle est l'espérance du temps jusqu'à la suivante~?
\end{interviewq}

\begin{interviewq}[$\star\star$ \iqroles{banque, trader}]\label{iq:qm:jump-processes:3}
Pourquoi une couverture en delta ne peut-elle pas éliminer le risque de saut~?
\end{interviewq}

\begin{interviewq}[$\star\star$ \iqroles{chercheur}]\label{iq:qm:jump-processes:4}
Énoncer la formule de Lévy--Khintchine et dire ce que représente chaque élément du \omterm{def:qm:jump-processes:triplet}{triplet de Lévy}.
\end{interviewq}

\begin{interviewq}[$\star\star$ \iqroles{chercheur, risque}]\label{iq:qm:jump-processes:5}
Comment le kurtosis des rendements évolue-t-il avec l'horizon dans un modèle de Lévy, et dans les données~?
\end{interviewq}

\begin{interviewq}[$\star\star\star$ \iqroles{banque, chercheur}]\label{iq:qm:jump-processes:6}
Dans une \omterm{def:qm:jump-processes:jd}{diffusion à sauts} pour une action sous la \omterm{def:qm:girsanov-and-changes-of-numeraire:emm}{probabilité risque-neutre}, que doit valoir la dérive, et pourquoi~?
\end{interviewq}
